% Generated by experiments/make_numbers.py -- do not edit.
%
% Tool audit.  AUCell's default ceiling is quoted from the released
% Bioconductor manual (AUCell_calcAUC, aucMaxRank = ceiling(0.05 *
% nrow(rankings))); UCell's from its documentation, as in the Results;
% the remaining rows describe the published algorithms only.

\begin{table}[t]
\centering
\small
\begin{tabular}{l>{\raggedright\arraybackslash}p{2.0cm}>{\raggedright\arraybackslash}p{3.1cm}>{\raggedright\arraybackslash}p{4.3cm}>{\raggedright\arraybackslash}p{2.5cm}}
\toprule
Tool & Ranking & Ceiling & Where the channel sits & Measured here \\
\midrule
AUCell \citep{aibar2017} & every gene of the panel, within the cell & released default $\lceil 0.05\,G \rceil$ & empty top ranks are filled by tie-break: the subject of this paper & yes --- the grids and the benchmark \\
UCell \citep{andreatta2021} & every gene, within the cell & ships 1{,}500 for 10x data; its guidance recommends the median detected & the same fill, at whatever ceiling the call passes & yes --- default against recommendation, per pair \\
singscore \citep{foroutan2018} & every gene, within the cell & no truncation & undetected genes tie at the bottom of the ranking; the fill does not arise, and the arbitrary order inside the tied block is a different exposure & no \\
GSVA \citep{hanzelmann2013} & every gene, kernel-weighted & no truncation & bottom-block ties as in singscore, smoothed by the kernel & no \\
module score \citep{tirosh2016} & none --- a mean of expression values & --- & no ranking; depth enters through normalisation & no \\
\bottomrule
\end{tabular}
\caption{Where the tie-break channel sits in the ranking tools the introduction names.  The \emph{Ceiling} column quotes each tool's own default where one exists: AUCell's from the released Bioconductor manual, checked against the current release; UCell's from its documentation, as in the Results.  Rows marked \emph{no} in the last column are statements about the published algorithms, read rather than measured; the table keeps that distinction visible instead of blurring it.}
\label{tab:tools}
\end{table}
